\chapter{Following a value through a program}
\MARSChapterHeading{Following a value}
\label{revision:lab-first-program}

This lab accompanies the main book's first program. Its purpose is to connect an
expression you understand with changes you can observe in registers. You will
predict the result of an instruction, compare the prediction with execution, and
use a disagreement to investigate a program. Familiarity with variables and
arithmetic is assumed; previous assembly experience is not.

The interactive lesson organizes the actions in the IDE. This chapter supplies
the reasoning behind those actions. You can also perform the activities manually
with the public example program.

In IntelliJ, open \emph{MARS Learning} from the \emph{Tools} menu and choose
\emph{Guided work}. Its activities connect your predictions with Run, Debug and
repair observations. \emph{Free writing} opens the related $3x+1$ problem near the
end of this chapter. Use it when you can trace the worked calculation and explain
which values must reach output. Switching modes preserves your source files;
guided work returns to the last guided segment you opened. If your course has not
released free writing, continue the guided activities.

\section{Prepare a prediction}

The expression is $2x+1$. Start with $x=20$ and distinguish the input, the doubled
value and the final result. They are 20, 40 and 41. These are three states in a
calculation, even if a high-level language lets you write the whole expression on
one line.

The example keeps the input in \texttt{\$v0} and the arithmetic result in
\texttt{\$t0}. A register is a named location, so a statement about a value needs a
location and an execution boundary. ``The value is 40'' is incomplete here.
``After the doubling instruction, \texttt{\$t0} contains 40'' is a prediction you
can check.

Before you run, write down what each of the two arithmetic instructions reads and
what it writes. Do not assume that a source is erased when read. When an instruction
uses one register as both a source and a destination, distinguish the old value
from the value written by that instruction.

In guided work, choose an answer directly in the \emph{Learn} pane. If the answer
is incorrect, read the explanation and try again. A border in the editor identifies
the instruction the question concerns; it does not execute that instruction or
establish a Debug pause. After a correct answer, the choices disappear and the
lesson explains the result before presenting the next question.

Figure~\ref{fig:revision-inline-prediction} shows the first question beside its
source. Keep the prediction and an execution observation separate: answering this
question does not run the program. The following activities test your reasoning
against a real execution.

% Native editable annotations over an unchanged light-theme IDE screenshot.
\begin{figure}[htbp]
  \centering
  \begin{tikzpicture}[alt={The Learn pane asks what addu t0, v0, v0 changes when v0 holds 20. Three clickable answers appear below the question. The editor borders that addu instruction, followed by addiu t0, t0, 1. This is source focus before execution, not a Debug pause.}]
    \node[anchor=north,inner sep=0] (quiz) at (0,0) {\MARSRevisionCrop{revision/captures/first-program/inline-prediction.png}{0.90\linewidth}{0.046875000}{0.524444444}{0.593750000}{0.197777778}};
    \begin{scope}[shift={(quiz.south west)},x={($(quiz.south east)-(quiz.south west)$)},y={($(quiz.north west)-(quiz.south west)$)}]
      \draw[black,line width=0.7pt] (0.076086957,0.056000000) rectangle (0.815217391,0.444000000);
      \node[draw=black,fill=white,circle,inner sep=1.5pt,font=\footnotesize] at ($(0.076086957,0.444000000)+(-4mm,0)$) {1};
    \end{scope}
    \node[anchor=north,inner sep=0] (source) at ([yshift=-7mm]quiz.south) {\MARSRevisionCrop{revision/captures/first-program/inline-prediction.png}{0.53\linewidth}{0.492187500}{0.517777778}{0.378906250}{0.431111111}};
    \begin{scope}[shift={(source.south west)},x={($(source.south east)-(source.south west)$)},y={($(source.north west)-(source.south west)$)}]
      \draw[black,line width=0.7pt] (0.060606061,0.478260870) rectangle (0.909090909,0.978260870);
      \node[draw=black,fill=white,circle,inner sep=1.5pt,font=\footnotesize] at ($(0.060606061,0.978260870)+(-4mm,0)$) {2};
    \end{scope}
  \end{tikzpicture}
  \par\smallskip
  \begin{minipage}{\linewidth}
    \small
    \begin{enumerate}
      \item Select an answer in the Learn pane. An incorrect choice produces feedback and leaves the question available for another attempt. A correct choice advances the lesson.
      \item The border identifies the instruction being discussed. Its first operand names the destination; the two remaining operands name its sources. This focus cue does not mean that execution is paused here.
    \end{enumerate}
  \end{minipage}
  \caption{A prediction stays beside its source. These native lesson controls are shown before execution; an answer is not a Run or Debug observation.}
  \label{fig:revision-inline-prediction}
\end{figure}


\section{Establish where execution is paused}

Run the example once to recognize its prompt and output. Then place a breakpoint
on the addition immediately below \texttt{calculate}, start Debug and enter 20.
The console request occurs before the breakpoint. When the debugger stops on the
addition, integer input has completed, but that addition has not executed yet.

This boundary matters. Reading the destination before an instruction and expecting
its new result would make a correct program appear wrong. The highlighted source
line tells you what will execute next. Some source lines represent several machine
instructions, so use the debugger's source step for this activity.

Find \texttt{\$v0} in the register view and interpret its value using the selected
number display. For input 20, decimal 20 and hexadecimal 14 name the same value.
A display format changes how bits are written on the screen; it does not change
the running program's bits.

If the input is not what you expected, investigate that before stepping through the
arithmetic. Confirm the input entered in this Debug session and the actual pause
location. A leftover console from an earlier run is not evidence about this one.

\section{Observe two different additions}

The first addition reads \texttt{\$v0} twice and writes \texttt{\$t0}. Predict both
registers after it executes. Step over the instruction and compare your prediction
with the view. If only the destination changed, explain why that agrees with the
instruction's operands.

The second addition reads \texttt{\$t0}, adds the immediate constant one, and writes
\texttt{\$t0}. Predict its result before stepping again. The same register name
appears twice because the old value supplies an input and the new value replaces
it. There is no separate hidden variable holding the whole expression.

Repeat the arithmetic trace with input $-3$. Its doubled value is $-6$ and its
final result is $-5$. A hexadecimal register display may show a long sequence
starting with \texttt{f}; use signed decimal or the main book's explanation of
representation when comparing it with your arithmetic. The instruction's name
does not turn this negative input into an unsigned magnitude.

Figure~\ref{fig:revision-first-program-boundaries} enlarges two successive pauses.
The source highlight and register values in each row belong to the same pause.
Use the comparison to separate a value already written from a value you are
predicting for the next instruction.

% Native editable figure. Coordinate receipts verify crops; do not regenerate edited prose.
\begin{figure}[htbp]
\centering
\begin{tikzpicture}[alt={Two Debug pauses for input 20. Before addu, v0 is hexadecimal 14 and t0 is zero. After one source step, addiu is the next instruction; v0 is still 14 and t0 is 28. These hexadecimal values mean decimal 20 and 40.}]
\node[anchor=north west,inner sep=0,font=\small\bfseries] (p0title) at (6mm,0) {Before the doubling instruction};
\node[anchor=north west,inner sep=0] (p0s) at ([yshift=-1mm]p0title.south west) {\MARSRevisionCrop{revision/captures/first-program/debug-before-add.png}{0.42\linewidth}{0.406250000}{0.514444444}{0.402343750}{0.414444444}};
\begin{scope}[shift={(p0s.south west)},x={($(p0s.south east)-(p0s.south west)$)},y={($(p0s.north west)-(p0s.south west)$)}]
\draw[black,line width=0.7pt] (0.146938776,0.390625000) rectangle (0.922448980,0.734375000);
\node[draw=black,fill=white,circle,inner sep=1.5pt,font=\footnotesize] at ($(0.146938776,0.734375000)+(-3mm,0)$) {1};
\end{scope}
\node[anchor=north west,inner sep=0] (p0v) at ([xshift=8mm]p0s.north east) {\MARSRevisionCrop{revision/captures/first-program/debug-before-add.png}{0.5\linewidth}{0.242187500}{0.075555556}{0.539062500}{0.764444444}};
\begin{scope}[shift={(p0v.south west)},x={($(p0v.south east)-(p0v.south west)$)},y={($(p0v.north west)-(p0v.south west)$)}]
\draw[black,line width=0.7pt] (0.042857143,0.847222222) rectangle (0.935714286,0.986111111);
\node[draw=black,fill=white,circle,inner sep=1.5pt,font=\footnotesize] at ($(0.042857143,0.986111111)+(-3mm,0)$) {2};
\draw[black,line width=0.7pt] (0.042857143,0.013888889) rectangle (0.935714286,0.152777778);
\node[draw=black,fill=white,circle,inner sep=1.5pt,font=\footnotesize] at ($(0.042857143,0.152777778)+(-3mm,0)$) {3};
\end{scope}
\node[anchor=north west,inner sep=0,font=\small\bfseries] (p1title) at ($(p0title.west |- p0v.south west)+(0,-4mm)$) {After one source step};
\node[anchor=north west,inner sep=0] (p1s) at ([yshift=-1mm]p1title.south west) {\MARSRevisionCrop{revision/captures/first-program/debug-after-add.png}{0.42\linewidth}{0.406250000}{0.514444444}{0.402343750}{0.414444444}};
\begin{scope}[shift={(p1s.south west)},x={($(p1s.south east)-(p1s.south west)$)},y={($(p1s.north west)-(p1s.south west)$)}]
\draw[black,line width=0.7pt] (0.146938776,0.046875000) rectangle (0.922448980,0.390625000);
\node[draw=black,fill=white,circle,inner sep=1.5pt,font=\footnotesize] at ($(0.146938776,0.390625000)+(-3mm,0)$) {4};
\end{scope}
\node[anchor=north west,inner sep=0] (p1v) at ([xshift=8mm]p1s.north east) {\MARSRevisionCrop{revision/captures/first-program/debug-after-add.png}{0.5\linewidth}{0.242187500}{0.075555556}{0.539062500}{0.764444444}};
\begin{scope}[shift={(p1v.south west)},x={($(p1v.south east)-(p1v.south west)$)},y={($(p1v.north west)-(p1v.south west)$)}]
\draw[black,line width=0.7pt] (0.042857143,0.847222222) rectangle (0.935714286,0.986111111);
\node[draw=black,fill=white,circle,inner sep=1.5pt,font=\footnotesize] at ($(0.042857143,0.986111111)+(-3mm,0)$) {5};
\draw[black,line width=0.7pt] (0.042857143,0.013888889) rectangle (0.935714286,0.152777778);
\node[draw=black,fill=white,circle,inner sep=1.5pt,font=\footnotesize] at ($(0.042857143,0.152777778)+(-3mm,0)$) {6};
\end{scope}
\end{tikzpicture}
\par\smallskip
\begin{minipage}{\linewidth}
\small
\begin{enumerate}
\item The highlighted \texttt{addu \$t0, \$v0, \$v0} is next. It has not written its destination yet.
\item \texttt{\$v0} contains \texttt{0x14}, decimal 20, returned by integer input. It supplies both source operands.
\item In this fresh session, \texttt{\$t0} is still zero. Reading it here does not reveal the pending addition\textquotesingle s result.
\item After Step Over, \texttt{addiu \$t0, \$t0, 1} is next. Its presence on the blue line does not mean the one has been added.
\item \texttt{\$v0} is still 20. Reading it twice did not consume it, and neither addition writes to this register.
\item \texttt{\$t0} now contains \texttt{0x28}, decimal 40. Predict its value after the next instruction before stepping again.
\end{enumerate}
\end{minipage}
\caption{A source pause is before the highlighted instruction. The source and register details in each row come from the same Debug pause.}
\label{fig:revision-first-program-boundaries}
\end{figure}

\clearpage

\section{Follow the answer into output}

Continue to the instructions below \texttt{report}. The program prints its text
prefix before printing the number. That requires changing \texttt{\$v0} to select
another service. The answer stays in \texttt{\$t0} until it is copied into
\texttt{\$a0} for integer output.

Pause just before the integer-output request. Identify the register that names the
service, the register containing its argument, and the register retaining the
calculation. For input 20, the service selector is one and the number to print is
41. Those values have different roles. Seeing a one in the selector does not mean
the calculation has become one.

Now continue and compare the console with the argument you inspected. This connects
the visible output to the interface that produced it. The print service does not
search for a register that happens to contain your intended answer.

Figure~\ref{fig:revision-first-program-output} shows the completed argument
preparation, just before the integer-output request. Explain the different roles
of the three marked registers before continuing.

% Native editable figure. Coordinate receipts verify crops; do not regenerate edited prose.
\begin{figure}[htbp]
  \centering
  \begin{tikzpicture}[alt={Correct output preparation at the pause before syscall. The preceding move copied t0 into a0, then li set v0 to one. The register view shows v0 one, a0 hexadecimal 29 and t0 hexadecimal 29. Hexadecimal 29 is decimal 41.}]
    \node[anchor=north west,inner sep=0,font=\small\bfseries] (p0title) at (6mm,0) {The correct request is ready};
    \node[anchor=north west,inner sep=0] (p0s) at ([yshift=-1mm]p0title.south west) {\MARSRevisionCrop{revision/captures/first-program/output-before-syscall.png}{0.42\linewidth}{0.406250000}{0.683333333}{0.402343750}{0.236666667}};
    \begin{scope}[shift={(p0s.south west)},x={($(p0s.south east)-(p0s.south west)$)},y={($(p0s.north west)-(p0s.south west)$)}]
      \draw[black,line width=0.7pt] (0.146938776,0.041666667) rectangle (0.922448980,0.958333333);
      \node[draw=black,fill=white,circle,inner sep=1.5pt,font=\footnotesize] at ($(0.146938776,0.958333333)+(-3mm,0)$) {1};
    \end{scope}
    \node[anchor=north west,inner sep=0] (p0v) at ([xshift=8mm]p0s.north east) {\MARSRevisionCrop{revision/captures/first-program/output-before-syscall.png}{0.5\linewidth}{0.242187500}{0.075555556}{0.539062500}{0.764444444}};
    \begin{scope}[shift={(p0v.south west)},x={($(p0v.south east)-(p0v.south west)$)},y={($(p0v.north west)-(p0v.south west)$)}]
      \draw[black,line width=0.7pt] (0.042857143,0.847222222) rectangle (0.935714286,0.986111111);
      \node[draw=black,fill=white,circle,inner sep=1.5pt,font=\footnotesize] at ($(0.042857143,0.986111111)+(-3mm,0)$) {2};
      \draw[black,line width=0.7pt] (0.042857143,0.013888889) rectangle (0.935714286,0.152777778);
      \node[draw=black,fill=white,circle,inner sep=1.5pt,font=\footnotesize] at ($(0.042857143,0.152777778)+(-3mm,0)$) {3};
      \draw[black,line width=0.7pt] (0.042857143,0.569444444) rectangle (0.935714286,0.708333333);
      \node[draw=black,fill=white,circle,inner sep=1.5pt,font=\footnotesize] at ($(0.042857143,0.708333333)+(-3mm,0)$) {4};
    \end{scope}
  \end{tikzpicture}
  \par\smallskip
  \begin{minipage}{\linewidth}
    \small
    \begin{enumerate}
      \item The \texttt{move} and \texttt{li} above the highlighted \texttt{syscall} have executed. The request itself is next.
      \item \texttt{\$v0} is one, selecting integer output. It no longer holds the input 20 or the earlier string-output selector four.
      \item \texttt{\$t0} retains \texttt{0x29}, decimal 41. Loading a new service selector did not change this arithmetic result.
      \item \texttt{\$a0} also contains 41. This is the argument the upcoming integer-output request will use; it does not search other registers for the answer.
    \end{enumerate}
  \end{minipage}
  \caption{Before integer output, the service selector and number to print have different roles. Enlarged details show one Debug pause for input 20.}
  \label{fig:revision-first-program-output}
\end{figure}

\clearpage

\section{Investigate a plausible wrong result}

Change only the instruction that prepares the integer-output argument so that it
copies \texttt{\$v0} instead of \texttt{\$t0}. Predict the console result for
input 20, then run the changed program. At this point, the prefix-printing request
has already reused \texttt{\$v0} as a service selector. Its earlier role as the
input register does not tell you its current contents.

A useful investigation separates computation from presentation. Check the input
at \texttt{calculate}, the answer after the two additions, and the argument before
integer output. If the answer is correct in \texttt{\$t0} but the argument is wrong
in \texttt{\$a0}, changing the arithmetic cannot repair the actual fault.

Restore the correct argument preparation and repeat the run. A check passed before
your edit is evidence about the earlier source. The interactive lesson must check
the edited source again before marking its execution task complete. Your earlier
predictions and any hints you opened can remain in your learning history.

Compare the faulty preparation in
Figure~\ref{fig:revision-first-program-wrong-argument} with the preceding correct
case. Find the first incorrect register value and its last writer. A correct value
elsewhere in the machine does not supply the argument automatically.

% Native editable figure. Coordinate receipts verify crops; do not regenerate edited prose.
\begin{figure}[htbp]
  \centering
  \begin{tikzpicture}[alt={Faulty output preparation at the pause before syscall. The preceding move reads v0 rather than t0. The register view shows t0 hexadecimal 29, decimal 41, v0 one and a0 four. The calculation and selector are correct, but the upcoming integer-output argument is wrong.}]
    \node[anchor=north west,inner sep=0,font=\small\bfseries] (p0title) at (6mm,0) {A correct calculation with the wrong argument};
    \node[anchor=north west,inner sep=0] (p0s) at ([yshift=-1mm]p0title.south west) {\MARSRevisionCrop{revision/captures/first-program/diagnosis-before-syscall.png}{0.42\linewidth}{0.406250000}{0.680000000}{0.402343750}{0.240000000}};
    \begin{scope}[shift={(p0s.south west)},x={($(p0s.south east)-(p0s.south west)$)},y={($(p0s.north west)-(p0s.south west)$)}]
      \draw[black,line width=0.7pt] (0.146938776,0.041666667) rectangle (0.922448980,0.958333333);
      \node[draw=black,fill=white,circle,inner sep=1.5pt,font=\footnotesize] at ($(0.146938776,0.958333333)+(-3mm,0)$) {1};
    \end{scope}
    \node[anchor=north west,inner sep=0] (p0v) at ([xshift=8mm]p0s.north east) {\MARSRevisionCrop{revision/captures/first-program/diagnosis-before-syscall.png}{0.5\linewidth}{0.242187500}{0.075555556}{0.539062500}{0.764444444}};
    \begin{scope}[shift={(p0v.south west)},x={($(p0v.south east)-(p0v.south west)$)},y={($(p0v.north west)-(p0v.south west)$)}]
      \draw[black,line width=0.7pt] (0.042857143,0.013888889) rectangle (0.935714286,0.152777778);
      \node[draw=black,fill=white,circle,inner sep=1.5pt,font=\footnotesize] at ($(0.042857143,0.152777778)+(-3mm,0)$) {2};
      \draw[black,line width=0.7pt] (0.042857143,0.847222222) rectangle (0.935714286,0.986111111);
      \node[draw=black,fill=white,circle,inner sep=1.5pt,font=\footnotesize] at ($(0.042857143,0.986111111)+(-3mm,0)$) {3};
      \draw[black,line width=0.7pt] (0.042857143,0.569444444) rectangle (0.935714286,0.708333333);
      \node[draw=black,fill=white,circle,inner sep=1.5pt,font=\footnotesize] at ($(0.042857143,0.708333333)+(-3mm,0)$) {4};
    \end{scope}
  \end{tikzpicture}
  \par\smallskip
  \begin{minipage}{\linewidth}
    \small
    \begin{enumerate}
      \item The faulty \texttt{move \$a0, \$v0} read the old selector four. The later \texttt{li \$v0, 1} selected integer output but did not rewrite \texttt{\$a0}.
      \item \texttt{\$t0} is still 41. The arithmetic already produced the intended result, so changing it would not repair this argument error.
      \item \texttt{\$v0} is one. The selected service is correct even though the number it will receive is wrong.
      \item \texttt{\$a0} is four. Trace its last writer back to the faulty \texttt{move}, then prepare it from the retained answer and check the edited program again.
    \end{enumerate}
  \end{minipage}
  \caption{The first wrong value is in the output argument. This pause is before the request; it does not show a completed console transcript.}
  \label{fig:revision-first-program-wrong-argument}
\end{figure}

\clearpage

\section{Work without the construction hints}

Choose \emph{Free writing} to open \texttt{independent.asm}. Keep the worked
example in its own file. Implement $3x+1$ while keeping the same input convention,
and change the answer message to describe the new calculation. Write a trace before running the program. Include positive, zero and negative inputs. Explain why your
intermediate results lead to the intended final value; three successful runs are
useful checks, but the explanation should describe the operation for the stated
input range.

Your course may require particular guided activities before this check becomes
available. The lesson lists those still needed. Complete each listed activity
with its current source, including its questions and Run or Debug observation.
Opening an activity alone does not complete it. Editing a required example,
counterexample or repair program requires a new execution check, even if you
restore its earlier text. Your instructor chooses which activities are required.

A course may limit completed independent checks. The results show how many you
have used. Each complete five-case assessment counts once, whether it passes or
fails; an assembly error or interrupted check does not. Editing or reopening the
file keeps the count. Use ordinary Run and Debug to investigate your work, and ask
your instructor for another assessment when needed.

Use \emph{Check my program} in the lesson to assess your independent file. The
results describe the public cases that passed or failed. Use \emph{Open a hint}
when you need help planning the next change.

The interactive exercise offers optional hints, from planning the value flow to
following a concrete trace. Reveal one when you have a specific question, such as
which value must survive the next instruction. After reading it, return to your
trace and explain what changed in your reasoning. Revealed hints remain in your
learning history after edits, but do not complete an execution check. Check the
edited program again and account for the values it produces.

Then choose a different temporary register and make the program work again. List
the uses that must change together. Register choice is an implementation decision;
the important interface is the number passed to output and the behavior promised
to the caller or user.

When you can make those changes and explain a wrong result by locating the first
incorrect state, you have used the debugger as a reasoning tool. Record any point
where you needed an explanation that the book or lab did not supply. Those gaps
are part of the pilot review, not something a successful execution check can hide.

Reopening an activity keeps your choices and revealed hints. To begin a new
progress history for this unit, choose \emph{Reset MARS Learning Progress} from
\emph{Tools} and read the confirmation before continuing. Your source files stay
in place. The reset saves one recovery copy; a later reset with new progress
replaces that copy. \emph{Recover MARS Learning Progress} restores the saved
choices and hints and keeps attempts made since the reset. Run the execution
checks again after recovery: a remembered answer does not establish that the
current source works.
